\documentclass[12pt,a4paper]{article}

\usepackage{amsmath, amssymb, amsthm}
\usepackage{geometry}
\usepackage{hyperref}
\usepackage{parskip}


\geometry{margin=2.5cm}

\newtheorem{theorem}{Theorem}
\newtheorem{lemma}{Lemma}
\newtheorem{corollary}{Corollary}
\theoremstyle{definition}
\newtheorem{definition}{Definition}
\newtheorem{remark}{Remark}

\title{\textbf{Les deux approches du LCT}
\author{Mathis Drieu}
\date{12/04/2026}

\begin{document}

\maketitle

\begin{table}[htb]
\centering
\footnotesize
\begin{tabularx}{\linewidth}{c c}
\toprule
Parameter & Value \\

\midrule
$c_S$ & 5 \\
$c_T$ & 1 \\
$c_{S,T}$ & cS \\
$c_{T,S}$ & cT \\

$g_1$ & 0.2  \\
$g_{1,S}$ & $g_1$ \\
$m_{T2}$ & 0.02  \\

$a_G$ & 1 \\
$a_S$ & 0.1 \\
$a_T$ & 2 \\

$e_G$ & 1 \\
$e_S$ & 0.2 \\
$e_T$ & 0.2 \\

$r$ & 1 \\

\bottomrule
\end{tabularx}
\caption{Vera application model parameters}
\end{table}

\section{Ando 2020}

Soit la fonction mémoire $M(t)$ :

\begin{equation}
    M(t) := \int_0^{+\infty} g\!\left(S(t-\theta),\, I(t-\theta)\right) K(\theta)\, d\theta,
    \tag{2.1}
\end{equation}

where $K : [0,+\infty) \to \mathbb{R}$ is a non-negative, normalized memory kernel:
\[
    \int_0^{+\infty} K(\theta)\, d\theta = 1.
\]

The \textbf{Erlang kernel} of order $n$ and rate $a > 0$ is defined as:

\begin{equation}
    \mathrm{Erl}_{n,a}(\theta) := \frac{a^n\, \theta^{n-1}}{(n-1)!}\, e^{-a\theta}, \quad \theta \geq 0.
    \label{eq:erlang}
\end{equation}

We take $K(\theta) = \mathrm{Erl}_{n,a}(\theta)$, so that:
\[
    M(t) = \int_0^{+\infty} g\!\left(S(t-\theta),\, I(t-\theta)\right) \mathrm{Erl}_{n,a}(\theta)\, d\theta.
\]

\begin{definition}[Auxiliary chain variables]
For $i = 1, \ldots, n$, define:
\begin{equation}
    Z_i(t) := \int_0^{+\infty} g\!\left(S(t-\theta),\, I(t-\theta)\right) \mathrm{Erl}_{i,a}(\theta)\, d\theta.
    \label{eq:Zi}
\end{equation}
\end{definition}

\begin{remark}
By definition, $Z_n(t) = M(t)$ for all $t$.
\end{remark}

\section{Key Lemmas}

\begin{lemma}[Boundary values of the Erlang kernel]
\label{lem:boundary}
\[
    \mathrm{Erl}_{i,a}(0) =
    \begin{cases}
        a & \text{if } i = 1, \\
        0 & \text{if } i \geq 2.
    \end{cases}
\]
\end{lemma}

\begin{proof}
From \eqref{eq:erlang}:
\[
    \mathrm{Erl}_{i,a}(0) = \frac{a^i \cdot 0^{i-1}}{(i-1)!} e^0.
\]
For $i = 1$: $\mathrm{Erl}_{1,a}(0) = a^1 \cdot 0^0 / 0! = a$.
For $i \geq 2$: $0^{i-1} = 0$, so $\mathrm{Erl}_{i,a}(0) = 0$.
\end{proof}

\begin{lemma}[Recurrence relation for the Erlang kernel]
\label{lem:recurrence}
For $i \geq 2$:
\[
    \frac{d}{d\theta}\mathrm{Erl}_{i,a}(\theta) = a\,\mathrm{Erl}_{i-1,a}(\theta) - a\,\mathrm{Erl}_{i,a}(\theta).
\]
For $i = 1$:
\[
    \frac{d}{d\theta}\mathrm{Erl}_{1,a}(\theta) = -a\,\mathrm{Erl}_{1,a}(\theta).
\]
\end{lemma}

\begin{proof}
Differentiating \eqref{eq:erlang} directly:
\begin{align*}
    \frac{d}{d\theta}\mathrm{Erl}_{i,a}(\theta)
    &= \frac{a^i}{(i-1)!}\left[(i-1)\theta^{i-2} - a\theta^{i-1}\right]e^{-a\theta} \\
    &= \frac{a^i(i-1)\theta^{i-2}}{(i-1)!}e^{-a\theta} - \frac{a^{i+1}\theta^{i-1}}{(i-1)!}e^{-a\theta} \\
    &= a\cdot\frac{a^{i-1}\theta^{i-2}}{(i-2)!}e^{-a\theta} - a\cdot\frac{a^i\theta^{i-1}}{(i-1)!}e^{-a\theta} \\
    &= a\,\mathrm{Erl}_{i-1,a}(\theta) - a\,\mathrm{Erl}_{i,a}(\theta), \quad i \geq 2.
\end{align*}
For $i=1$: $\mathrm{Erl}_{1,a}(\theta) = ae^{-a\theta}$, so $\frac{d}{d\theta}\mathrm{Erl}_{1,a}(\theta) = -a^2 e^{-a\theta} = -a\,\mathrm{Erl}_{1,a}(\theta)$.
\end{proof}

\section{Main Theorem: The Linear Chain Trick}

\begin{theorem}[Linear Chain Trick]
\label{thm:LCT}
Let $Z_i(t)$ be defined as in \eqref{eq:Zi}. Then the $Z_i$ satisfy the following
system of ODEs:
\begin{equation}
    \begin{cases}
        Z_1'(t) = a\, g(S(t), I(t)) - a\, Z_1(t), \\[6pt]
        Z_i'(t) = a\, Z_{i-1}(t) - a\, Z_i(t), \quad i = 2, \ldots, n,
    \end{cases}
    \label{eq:chain}
\end{equation}
with $Z_n(t) = M(t)$.
\end{theorem}

\begin{proof}
We differentiate $Z_i(t)$ using the substitution $u = t - \theta$ (so $\theta = t - u$,
$d\theta = -du$), which rewrites \eqref{eq:Zi} as:
\[
    Z_i(t) = \int_{-\infty}^{t} g(S(u), I(u))\, \mathrm{Erl}_{i,a}(t - u)\, du.
\]

Applying the Leibniz differentiation rule:
\begin{equation}
    Z_i'(t) = g(S(t), I(t))\cdot \mathrm{Erl}_{i,a}(0)
    + \int_{-\infty}^{t} g(S(u), I(u))\, \frac{\partial}{\partial t}\mathrm{Erl}_{i,a}(t-u)\, du.
    \label{eq:leibniz}
\end{equation}

\medskip
\noindent\textbf{Case $i \geq 2$.}

By Lemma~\ref{lem:boundary}, $\mathrm{Erl}_{i,a}(0) = 0$, so the boundary term vanishes.
Since $\frac{\partial}{\partial t}\mathrm{Erl}_{i,a}(t-u) = \frac{d}{d\theta}\mathrm{Erl}_{i,a}(\theta)\big|_{\theta=t-u}$,
applying Lemma~\ref{lem:recurrence}:

\begin{align*}
    Z_i'(t) &= \int_{-\infty}^{t} g(S(u), I(u))
    \left[a\,\mathrm{Erl}_{i-1,a}(t-u) - a\,\mathrm{Erl}_{i,a}(t-u)\right] du \\
    &= a\int_{-\infty}^{t} g(S(u), I(u))\,\mathrm{Erl}_{i-1,a}(t-u)\,du
    - a\int_{-\infty}^{t} g(S(u), I(u))\,\mathrm{Erl}_{i,a}(t-u)\,du \\
    &= a\, Z_{i-1}(t) - a\, Z_i(t).
\end{align*}

\medskip
\noindent\textbf{Case $i = 1$.}

By Lemma~\ref{lem:boundary}, $\mathrm{Erl}_{1,a}(0) = a$, so the boundary term contributes
$a\, g(S(t), I(t))$. Applying Lemma~\ref{lem:recurrence} for $i=1$:

\begin{align*}
    Z_1'(t) &= a\, g(S(t), I(t))
    + \int_{-\infty}^{t} g(S(u), I(u))\left[-a\,\mathrm{Erl}_{1,a}(t-u)\right] du \\
    &= a\, g(S(t), I(t)) - a\, Z_1(t).
\end{align*}

\medskip
Finally, $Z_n(t) = M(t)$ by definition, which completes the proof.
\end{proof}

\section{Interpretation}

The chain \eqref{eq:chain} has a clear structure:

\begin{itemize}
    \item \textbf{$Z_1$} is pulled directly toward the current nonlinear input $g(S(t), I(t))$
          with rate $a$. It is the least smoothed variable.
    \item \textbf{$Z_i$} ($i \geq 2$) is pulled toward the previous link $Z_{i-1}$, adding
          one further layer of smoothing at each step.
    \item \textbf{$Z_n = M(t)$} is the most smoothed and most delayed variable, representing
          the full distributed memory, and is the quantity that feeds back into the main
          dynamical equations.
\end{itemize}

The information flow is therefore:
\[
    g(S,I) \;\longrightarrow\; Z_1 \;\longrightarrow\; Z_2 \;\longrightarrow\; \cdots
    \;\longrightarrow\; Z_n = M(t).
\]

\begin{remark}[Mean delay]
The Erlang kernel $\mathrm{Erl}_{n,a}$ has mean $n/a$ and variance $n/a^2$. As $n \to \infty$
with $n/a$ fixed, it converges to a Dirac mass at the mean delay, recovering a discrete
time lag. The chain thus interpolates continuously between a pure exponential memory
($n=1$) and a discrete delay ($n \to \infty$).
\end{remark}

\begin{remark}[Linearity of the chain]
All equations in \eqref{eq:chain} except the first are \emph{linear} ODEs. The nonlinearity
of $g(S,I)$ enters only through the forcing term of $Z_1'$. This is precisely what
McDonald (1978) calls the ``non-linear'' total lag case, in contrast to diagonal or
off-diagonal lags where only one variable is convolved and the chain is entirely linear.
\end{remark}

\end{document}